\section{Momentum, energy and admissible physical coupling}
\label{sec:r13-energy}
The exact seed needs a rule for composing physical models without creating
energy at their interfaces. R13 adopts established momentum and thermodynamic
balances, then derives a finite energy-port contract compatible with R12's
mechanical and thermal states. These are restrictions on a selected physical
model. They do not assign physical meaning to a bit pattern by themselves.

\subsection{Momentum uses a realized frame and a declared body}
For a closed material body in classical continuum mechanics, let $\rho>0$
be mass density, $v$ velocity, $b$ body force per unit mass, and $\sigma$
Cauchy stress. In an inertial Cartesian frame the local balance laws are
\begin{equation}
 \dot\rho+\rho\nabla\!\cdot v=0,\qquad
 \rho\dot v=\nabla\!\cdot\sigma+\rho b,\qquad \sigma=\sigma^T.
 \label{eq:r13-momentum}
\end{equation}
Dots here are material derivatives. Stress symmetry assumes the ordinary
continuum without independent body couples or couple stresses. A moving
control volume instead requires the corresponding mass and momentum fluxes.
These balances and their energy/entropy counterparts are stated in
\href{https://www.eng.utah.edu/~banerjee/Notes/ThermoElastic.pdf}
{EN1: Banerjee, \emph{Basic Thermoelasticity}, sections 1.1--1.3}.

For a rigid body of fixed mass $m$, centre of mass $r_C$, body-to-inertial
rotation $R_B\in SO(3)$ and constant body-frame inertia $I_B\succ0$,
the corresponding finite specialization is
\begin{align}
 m\ddot r_C&=F_{\rm ext},&
 I_B\dot\omega_B+\omega_B\times(I_B\omega_B)&=\tau_B,\\
 \dot R_B&=R_B[\omega_B]_\times,&
 [\omega_B]_\times a&=\omega_B\times a .
 \label{eq:r13-euler}
\end{align}
The torque is about the centre of mass; body components and inertial
components must not be mixed. The cross term follows from differentiating
angular momentum in rotating axes. See
\href{https://ocw.mit.edu/courses/16-07-dynamics-fall-2009/5e1d8699338146e5127080b880b906d6_MIT16_07F09_Lec28.pdf}
{EN2: MIT Dynamics, lecture 28, Euler's equations}.
For this specialization direct contraction gives
\begin{equation}
 \frac{d}{dt}\left(\tfrac12m|\dot r_C|^2+
                 \tfrac12\omega_B^TI_B\omega_B\right)
   =F_{\rm ext}\!\cdot\dot r_C+\tau_B\!\cdot\omega_B,
 \label{eq:r13-rigid-power}
\end{equation}
because $\omega_B\cdot[\omega_B\times(I_B\omega_B)]=0$.
An inertia estimate, applied torque and orientation measurement are different
inputs; a Madgwick attitude estimate alone supplies no torque law.

R12's frame annotations preserve coordinate identity. Their physical
realization also requires origins, alignment and time convention. Transforming
to accelerating or rotating coordinates changes the acceleration equation;
renaming an inertial force law as a station-frame law is insufficient.
Likewise a superposed rigid rotation must not generate elastic strain energy.
For a finite-deformation hyperelastic model with deformation gradient
$F=\nabla_X x$, $\det F>0$, the selected objective free energy obeys
\begin{equation}
 \Psi(QF,T)=\Psi(F,T),\qquad Q\in SO(3),
 \qquad (QF)^T(QF)=F^TF .
 \label{eq:r13-objective}
\end{equation}
Material directions, if present, remain explicit reference data; objectivity
does not imply isotropy. This condition and the two-point transformation of
$F$ are discussed in
\href{https://www.weizmann.ac.il/chemphys/bouchbinder/sites/chemphys.bouchbinder/files/uploads/Courses/2019/Lectures/ContinuumPhysics_1-10_2019.pdf}
{EN3: Bouchbinder, \emph{Non-Equilibrium Continuum Physics}, section 4}.
The existing $\epsilon=\operatorname{sym}\nabla u$ model remains a small-gradient
approximation; equation \eqref{eq:r13-objective} does not extend it to arbitrary
finite rotations. A finite-rotation specialization must supply its own
objective constitutive equation and compatible finite state.

\subsection{A finite port contract with an explicit power proof}
Use the finite state $x$, differentiable storage $H(x)$, and matrices
$\mathcal J(x)=-\mathcal J(x)^T$, $\mathcal R(x)=\mathcal R(x)^T\succeq0$
and $\mathcal G(x)$. Define
\begin{equation}
 \dot x=(\mathcal J-\mathcal R)\nabla H+\mathcal G u,
 \qquad y=\mathcal G^T\nabla H .
 \label{eq:r13-port}
\end{equation}
This is the standard finite port-Hamiltonian form; $u^Ty$ is power supplied
at the declared boundary. See
\href{https://arxiv.org/html/2412.19673v1#S1.SS1}
{EN4: van der Schaft, \emph{Port-Hamiltonian nonlinear systems}, equations 5--6}.
The symbol $\mathcal J$ is unrelated to the JK input word $J$.
Each contraction must be dimensionally compatible. If components of $x$ have
different units, their derivative and matrix entries have corresponding
component units; they are stored as typed tuples and explicit contractions,
not forced into a homogeneous vector type.

\begin{proposition}[Finite power balance and composition]
For a differentiable solution of \eqref{eq:r13-port},
\begin{equation}
 \dot H=u^Ty-\mathcal D,\qquad
 \mathcal D=(\nabla H)^T\mathcal R\nabla H\ge0 .
 \label{eq:r13-power}
\end{equation}
For two such systems whose ports admit a dimensionally compatible map $L$,
connect $u_1=Ly_2$, $u_2=-L^Ty_1$. Their internal power cancels.
Their combined storage is
$H_1+H_2$ and its rate is minus their dissipation plus any remaining external
port power.
\end{proposition}
\begin{proof}
Apply the chain rule and use $z^T\mathcal Jz=0$ for a skew matrix.
The two connected powers sum to $y_2^TL^Ty_1-y_1^TLy_2=0$.
The entries of $L$ carry any required unit and
frame conversion; a force is not silently identified with a velocity.
\end{proof}
If $H$ has a lower bound, integration gives the usual passivity inequality
$H(t)-H(0)\le\int_0^t u^Ty\,ds$.
This limits extraction relative to storage and supply; it does not by itself
prove asymptotic stability, accurate tracking, or a sufficient actuator.
Algebraic port connections must also yield a well-defined finite evolution.

Physical port pairs include force/velocity, torque/angular velocity and
voltage/current. In a thermal port the energy pairing is absolute temperature
times entropy flow, $T\dot S_{\rm transfer}=\dot Q$ at that port. Pairing
temperature directly with heat power and calling the product power would
have the wrong units. Entropy production inside the body is distinct from
entropy transported through a port.

\subsection{The existing modal mechanics fits this contract}
R12 already declares a finite displacement expansion
$u_m(X,t)=\sum_iq_i(t)\phi_i(X)$, mass and stiffness matrices from its weak
form, and support conditions. For a fixed basis, fixed material parameters,
isothermal or explicitly uncoupled mechanics, take
\begin{equation}
 M\ddot q+D\dot q+Kq=B u,\qquad
 M=M^T\succ0,\quad K=K^T\succeq0,\quad D=D^T\succeq0.
 \label{eq:r13-modal}
\end{equation}
Let $p=M\dot q$ and
\begin{equation}
 H_m(q,p)=\tfrac12p^TM^{-1}p+\tfrac12q^TKq,
 \quad
 \mathcal J=\begin{bmatrix}0&I\\-I&0\end{bmatrix},\quad
 \mathcal R=\begin{bmatrix}0&0\\0&D\end{bmatrix},\quad
 \mathcal G=\begin{bmatrix}0\\B\end{bmatrix}.
\end{equation}
Then \eqref{eq:r13-port} gives exactly \eqref{eq:r13-modal}, with
\begin{equation}
 y=B^T\dot q,\qquad
 \dot H_m=u^TB^T\dot q-\dot q^TD\dot q .
 \label{eq:r13-modal-power}
\end{equation}
The PSD stiffness admits free modes; asymptotic displacement stability does
not follow when rigid or undamped modes remain. If $M$ depends on configuration,
the derivative of its kinetic energy is part of the dynamics; merely replacing
constant $M$ by $M(q)$ in \eqref{eq:r13-modal} generally misses terms.

Finite modal reduction preserves this displayed energy identity when its
coefficient assumptions hold. It does not prove the basis resolves every
physical mode. Spatial integrals still require the charts, measures and
integrability obligations already specified by R12. A continuum PDE printed
above is an origin for the model; it is not an implemented infinite-dimensional
XOP operator.

An actual ASA/NA--JK specialization provides an explicit decode from its
completed old state to a command, then a transduction law to the physical
effort $u(t)$. A voltage command cannot be substituted as a force without
that law. The resulting work is
\begin{equation}
 W_n=\int_{t_n}^{t_{n+1}}u(t)^Ty(t)\,dt .
 \label{eq:r13-command-work}
\end{equation}
The command may be held over the interval; the velocity generally is not.
A bit selecting a positive effort can supply energy through the actuator's
power source. Neither bit packing nor a delayed update makes that port passive.

\subsection{Convert dissipated work to heat exactly once}
For the uncoupled mechanical storage above let
$P_d=\dot q^TD\dot q\ge0$. Suppose all this dissipated mechanical power
thermalizes in the modeled assembly, and fixed or state-declared fractions
$w_i\ge0$, $\sum_iw_i=1$, distribute it among thermal nodes. Write
\begin{align}
 \dot H_m&=P_{\rm mech,in}-P_d,\\
 C_i\dot T_i&=P_{i,\rm other}+w_iP_d
       +\sum_{j\ne i}G_{ij}(T_j-T_i)+G_{i0}(T_a-T_i),
 \label{eq:r13-heat-transfer}
\end{align}
where $C_i>0$, $G_{ij}=G_{ji}\ge0$ and $G_{i0}\ge0$.
The stores $U_i=C_i(T_i-T_{i,\rm ref})$ are an uncoupled constant-capacity
model. Their sum with $H_m$ obeys the directly derived identity
\begin{equation}
 \frac{d}{dt}(H_m+\sum_iU_i)
 =P_{\rm mech,in}+\sum_iP_{i,\rm other}
                 +\sum_iG_{i0}(T_a-T_i).
 \label{eq:r13-total-transfer}
\end{equation}
Internal conduction cancels in pairs, as does mechanical-to-thermal power.
Thus $P_d$ is not added again as an external input. If only a fraction
thermalizes locally, the remaining energy needs its destination, such as
outgoing sound or a separate store. A resistor's $VI$ and an optical absorber's
thermalized energy are handled with the same nonduplicating boundary ledger.
The distinction between work conversion and entropy creation is treated in
\href{https://live.ocw.mit.edu/courses/2-141-modeling-and-simulation-of-dynamic-systems-fall-2006/1e6c6f7dd1e3033e1af593d1d28c1d6f_work_to_heat_tra.pdf}
{EN5: MIT, \emph{Work-to-heat Transduction}}.

\subsection{Entropy supplies a second, independent condition}
Conserving total energy does not make every proposed heat exchange physical.
Restrict the node model to local equilibrium with $T_i>0$, and define
\begin{equation}
 S_i-S_{i,\rm ref}=C_i\ln(T_i/T_{i,\rm ref}),\qquad
 T_{i,\rm ref}>0,\qquad \dot S_i=C_i\dot T_i/T_i.
\end{equation}
For one reciprocal conductive link the derived total entropy rate is
\begin{equation}
 \dot S_{ij}=\frac{G_{ij}(T_j-T_i)}{T_i}
          +\frac{G_{ij}(T_i-T_j)}{T_j}
       =\frac{G_{ij}(T_i-T_j)^2}{T_iT_j}\ge0 .
 \label{eq:r13-node-entropy}
\end{equation}
A link to an ideal bath at $T_a>0$ has the same expression after including
the bath's entropy change. Local thermalization of $w_iP_d$ contributes
$w_iP_d/T_i\ge0$. Arbitrary external heat, light, matter and work ports must
carry their own energy and entropy accounting; an incident optical field's
entropy flux is not assigned the thermal formula $P_{\rm optical}/T_i$ merely
because its energy is later absorbed.

In local continuum heat conduction, with heat flux $q_h=-\kappa\nabla T$
and positive-semidefinite symmetric conductivity $\kappa$, the analogous
production density is
\begin{equation}
 \xi_{\rm cond}=q_h\!\cdot\nabla(1/T)
          =\frac{(\nabla T)^T\kappa\nabla T}{T^2}\ge0.
 \label{eq:r13-fourier}
\end{equation}
Fourier's constitutive law is a local-equilibrium model, not a universal
description of ballistic transport or ultrafast material response. Entropy
flux and production are distinguished in
\href{https://live.ocw.mit.edu/courses/2-141-modeling-and-simulation-of-dynamic-systems-fall-2006/19a8fbec0c3f9d83e01f6d6fe759bcc1_entropy_producti.pdf}
{EN6: MIT, \emph{Entropy Production and Nonlinearity}} and
\href{https://ocw.mit.edu/courses/2-43-advanced-thermodynamics-spring-2024/1hMAkYNBhBaLcLpX6vPnQgR65rSvPUs5c_transcript.pdf}
{EN7: MIT Advanced Thermodynamics, lecture 21, heat conduction}.

\subsection{Thermoelastic expansion requires reciprocal heat coupling}
R12's thermal strain is useful, but a coupled model cannot independently
choose an elastic energy, expansion force and unrelated heat capacity.
The following small-gradient, local-equilibrium specialization derives them
from one Helmholtz free-energy density $\psi(\epsilon,T)$ per reference volume.
All densities below use that same small-strain reference convention. Define
\begin{equation}
 s=-\psi_T,\qquad e=\psi+Ts=\psi-T\psi_T,\qquad
 \sigma=\psi_\epsilon+\sigma_d,\qquad
 c_\epsilon=-T\psi_{TT}>0.
 \label{eq:r13-free-energy}
\end{equation}
Here $s$ is entropy density, $e$ internal energy density, $\sigma_d$ dissipative
stress, and $c_\epsilon$ volumetric heat capacity at fixed strain.
These thermodynamic conjugacies follow the free-energy formulation in
\href{https://www.eng.utah.edu/~banerjee/Notes/ThermoElastic.pdf}
{EN1, sections 1.3.3--1.3.4}; the following compatibility calculation is
the R13 specialization.

Combining the local balances
$\dot e=\sigma:\dot\epsilon-\nabla\!\cdot q_h+r$ and
$\dot s+\nabla\!\cdot(q_h/T)-r/T=\xi\ge0$ gives
\begin{align}
 T\xi&=\sigma:\dot\epsilon-\dot\psi-s\dot T
                       -q_h\!\cdot\nabla T/T,\\
 &=\sigma_d:\dot\epsilon-q_h\!\cdot\nabla T/T\ge0,\\
 c_\epsilon\dot T&=T\psi_{T\epsilon}:\dot\epsilon
          +\sigma_d:\dot\epsilon-\nabla\!\cdot q_h+r.
 \label{eq:r13-reciprocal-heat}
\end{align}
The last equation follows by differentiating
$e=\psi-T\psi_T$ and canceling $\psi_\epsilon:\dot\epsilon$.
Here $r$ is heat supply after any declared local conversion, rather than a
second copy of mechanical dissipation. With
$\sigma_d=\mathbb V:\dot\epsilon$, symmetric $\mathbb V\succeq0$ on
symmetric strains, and Fourier conductivity as above, both displayed
dissipative contributions are nonnegative. The reversible term
$T\psi_{T\epsilon}:\dot\epsilon$ need not be positive: it exchanges thermal
and mechanical storage without being dissipation.

A concrete minimal profile uses constant elastic tensor $\mathbb C$ positive
on admissible strains, constant symmetric thermal-stress tensor $\beta$,
constant $c_v>0$ and $T_0>0$:
\begin{align}
 \psi(\epsilon,T)&=\tfrac12\epsilon:\mathbb C:\epsilon
    -(T-T_0)\beta:\epsilon
    +c_v\bigl[T-T_0-T\ln(T/T_0)\bigr],\\
 \sigma&=\mathbb C:\epsilon-(T-T_0)\beta+\sigma_d,\\
 s&=\beta:\epsilon+c_v\ln(T/T_0),\\
 e&=\tfrac12\epsilon:\mathbb C:\epsilon
           +T_0\beta:\epsilon+c_v(T-T_0),\qquad c_\epsilon=c_v,\\
 c_v\dot T&=-T\beta:\dot\epsilon
          +\sigma_d:\dot\epsilon-\nabla\!\cdot q_h+r.
 \label{eq:r13-minimal-thermoelastic}
\end{align}
Choosing $\beta=\mathbb C:(\alpha_T I)$ reproduces R12's constant-coefficient
thermal-expansion stress. The linear term in $e$ is required by the selected
entropy convention and coupling; replacing $e$ by an unrelated
``elastic plus $c_vT$'' formula breaks the same first-law calculation.
The profile applies on a declared strain and positive-temperature range;
it is not a fitted constitutive law for an unspecified specimen.

Additional reporter, damage or chemical state $z$ changes this closure to
$\psi(\epsilon,T,z)$ and adds $-\psi_z\!\cdot\dot z$ to $T\xi$.
An interval-preserving reporter update from R12 proves its range, not this
new thermodynamic inequality. Either specify the driving energy/rates that
meet it, or keep the reporter as an empirical readout state with its energy
exchange separately declared. No new chemical model is assumed here.

\subsection{Switching and exact steps keep their own energy obligations}
If a physical mode $a$ changes the storage function, continuous parameter
motion has the additional power term $H_a\dot a$ in the chain rule. At an
idealized discrete switch, use one consistent reference for the stores and
record
\begin{equation}
 \Delta H_{\rm switch}=H_{a^+}(x^+)-H_{a^-}(x^-)
       =W_{\rm switch,in}-E_{\rm switch,out}.
 \label{eq:r13-switch}
\end{equation}
An ideal routing switch can have zero storage jump. A change of stiffness at
fixed $q,p$ and fixed $M$ instead contributes
$\tfrac12q^T(K^+-K^-)q$; a state reset or impulsive motion has its own
mechanical balance. Labeling the selector as ASA, NA, JK or a thermal bit
does not remove these terms. Passive-mode energy inequalities cannot simply
be concatenated across unaccounted jumps.

Exact evaluation preserves the expression of a selected time-step map.
It does not automatically preserve \eqref{eq:r13-power} between its endpoints.
Section~\ref{sec:r13-step} supplies a finite midpoint specialization with its
own derived step energy identity for fixed linear mechanical coefficients.
The continuous coupled thermoelastic equations here retain their separate
finite-reduction, existence and time-discretization obligations.

\subsection{What is added to the kernel contract}
No opcode is needed for these laws. R13 binds energy and entropy functions,
their derivatives, domain conditions, port pairings and selected finite
evolution as versioned templates using existing exact arithmetic,
\code{LN}, finite integrals and finite \code{ODE} declarations.
States of unlike units use tuples; matrix contractions expand componentwise.
A valid specialization can now present its power identity, heat-transfer
cancellation and nonnegative entropy production as explicit mathematical
obligations. A false sign, unpowered actuator, duplicated heat term or
incompatible thermal expansion can be identified before a kernel is run.
This is the useful added foundation. Specimen parameters, model adequacy,
runtime evaluation and physical performance remain separate evidence.
